% !TEX root = ../main.tex
\section{등록 분석}
\label{sec:map-registered}

등록 분석(미리 등록한 그대로 한 분석)은 분석 계획과 W1 등록 문서가 정해 둔 분석이에요. 이 연구가 판정에 쓸 자료를 뽑기 전에 정했어요(부록~\ref{sec:pipeline-stages}). 이 분석은 금액을 토큰의 원래 단위(소수점 조정을 하지 않은 가장 작은 단위)로 재고, 모든 ERC-20(토큰 공통 규칙) 토큰을 쓰고, 값 변환(금액이 아무리 커도 1을 넘지 않게 바꾸는 식)에 $b=1$을 써요. 또 2023년 10월 1일부터의 모든 이벤트(계약이 남기는 기록)로 점수를 매기고, 계약 $27{,}844$개로 된 계획의 코호트(묶음)를 써요. 수정 분석(등록 뒤에 금액을 달러로 바꿔 다시 한 분석)은 금액을 달러(USD)로 매기고, 다섯 토큰을 쓰고, 동결일(점수를 고정하는 날)마다 그 앞 열두 달의 기록으로 점수를 매겨요. 수정 분석은 등록 분석의 결과를 안 뒤에 정했어요(\ref{sec:amount-revision}절). 그리고 본문의 장들은 수정 분석을 보고해요. 표~\ref{tab:registered-vs-revised}는 등록 분석의 판정 대비와 판정을, 고친 점수에서 구한 같은 대비 옆에 나란히 놓아요. 이 표의 마지막 네 행은 라벨(나중에 채점에 쓰는 정답 값)을 본 뒤에 계산한 차수 대비이고, 판정이 없어요. 두 열은 R1을 빼면 같은 판정에 이르러요. R1은 고친 점수에서는 성립하고, 등록 분석에서는 성립하지 않아요. 이 연구가 판정에 쓸 자료를 뽑기 전에 정한 시험의 결과는 등록 열의 판정뿐이에요. 모든 계수, 구간, 판정을 담은 등록 분석의 전체 결과는 바꾸지 않은 그대로 \texttt{data/3-published-results-for-thesis/}에 공개했어요. 그 요약 파일인 \texttt{results\_digest.md}도 함께 공개했어요.

% Generated by experiments/scripts/export_latex.py -- do not edit by hand.
% Source: experiments/data/2-processed-tables-and-evaluations/same-window/eval_summary.json
% Source: experiments/data/2-processed-tables-and-evaluations/holdout-w0/eval_summary.json
% Source: experiments/data/2-processed-tables-and-evaluations/registered-w1/eval_summary.json
% Source: experiments/data/2-processed-tables-and-evaluations/revised-usd-12m/same-window/eval_summary.json
% Source: experiments/data/2-processed-tables-and-evaluations/revised-usd-12m/holdout-w0/eval_summary.json
% Source: experiments/data/2-processed-tables-and-evaluations/revised-usd-12m/window-w1/eval_summary.json
\begin{table}[htbp]
\centering
\caption[등록 분석과 수정 분석의 판정 대비.]{등록 분석의 판정 대비와 판정이에요(금액은 원래 단위, 모든 ERC-20 토큰, $b=1$, 점수는 2023년 10월 1일부터의 모든 이벤트로 매김). 분석 계획과 W1 등록 문서에 따라 얻은 그대로예요. 그 옆에 고친 점수에서 구한 같은 대비를 놓았어요(금액은 달러(USD), 목록에 오른 다섯 토큰, 점수는 각 동결일 전 열두 달의 이벤트로 매김). 수정 분석은 점수 기간까지 포함해서, 등록 분석의 결과를 안 뒤에 정했어요. 수정 분석 열은 등록한 규칙과 기준값을 고친 점수에 적용한 것이고, 등록한 검사가 아니에요. 마지막 묶음에는 두 분석의 어떤 규칙도 판정에 쓰지 않는 차수 대비가 있어요. 이 대비들은 두 분석 모두에서 보고해요. $\Delta\tau$와 95\% 짝지은 부트스트랩 구간을 적었어요(다시 뽑기 400번과 2{,}000번, 시드 42). 규칙과 그 기준값은 \ref{sec:holdout-protocol}절과 \ref{sub:fresh-holdout}절에 있어요.}
\label{tab:registered-vs-revised}
\fitwidth{%
\begin{tabular}{lcc}
\toprule
항목 & 등록 분석 & 수정 분석 \\
\midrule
매칭 코호트 (계약 수) & 27{,}844 & 6{,}483 \\
W0 스펜더 코호트 & 33{,}101 & 7{,}914 \\
W1 스펜더 코호트 & 34{,}587 & 7{,}960 \\
W1 트레이더 & 70 & 70 \\
EndorseRank 대 AWP, 같은 기간 ($\tau$) & 0.217 [0.209, 0.225] & 0.199 [0.181, 0.221] \\
\midrule
A, W0: C-PR 빼기 C-PR ($\lambda=1$), 새 허락 쌍 & $+0.008$ [0.003, 0.013] & $+0.031$ [0.021, 0.042] \\
A, W0: C-PR 빼기 C-PR ($\lambda=0$), 새 송금자 & $-0.030$ [-0.034, -0.026] & $-0.030$ [-0.040, -0.020] \\
규칙 A, 주된 부분 & 충족 안 됨 & 충족 안 됨 \\
규칙 A, 보조 부분 & 충족 안 됨 & 충족 안 됨 \\
\midrule
B, F1: C-PR 빼기 C-PR ($\lambda=0$), 새 허락 쌍 & $+0.101$ [0.092, 0.111] & $+0.170$ [0.149, 0.189] \\
B, F2: C-PR 빼기 C-PR ($\lambda=0$), 새 송금자 & $-0.025$ [-0.029, -0.021] & $-0.025$ [-0.038, -0.015] \\
B, F3: C-PR 빼기 C-PR ($\lambda=0$), 청산 없이 닫은 비율 & $-0.044$ [-0.100, -0.002] & $+0.000$ [-0.010, 0.010] \\
규칙 B (F1, F2, F3) & 충족 안 됨 & 충족 안 됨 \\
규칙 B, 민감도 1 (AWP를 비교 대상으로 함) & 충족 안 됨 & 충족 안 됨 \\
규칙 B, 민감도 2 (외톨이 점) & 충족 안 됨 & 충족 안 됨 \\
F6 (규칙 A, 주된 부분, W1에서) & 불성립 & 불성립 \\
\midrule
P, W1: 받은 허락 수 빼기 송금 입차수 & $+0.063$ [-0.026, 0.148] & $+0.141$ [0.037, 0.240] \\
Q, W1: EndorseRank (활동 재시작) 빼기 AWP & $+0.051$ [-0.085, 0.213] & $+0.023$ [-0.062, 0.113] \\
R1 (그냥 센 수의 수준에서 앞섬) & 충족 안 됨 & 충족됨 \\
R2 (페이지랭크 수준의 앞섬) & 충족 안 됨 & 충족 안 됨 \\
R3 (특이성) & 적용됨 & 적용됨 \\
\midrule
W0: EndorseRank 빼기 받은 허락 수, 새 쌍 & $-0.011$ [-0.016, -0.006] & $-0.090$ [-0.100, -0.079] \\
W0: AWP 빼기 송금 입차수, 새 송금자 & $+0.038$ [0.035, 0.042] & $-0.038$ [-0.043, -0.032] \\
W1: EndorseRank 빼기 받은 허락 수, 새 쌍 & $-0.029$ [-0.032, -0.025] & $-0.082$ [-0.090, -0.073] \\
W1: AWP 빼기 송금 입차수, 새 송금자 & $+0.020$ [0.016, 0.023] & $-0.041$ [-0.047, -0.034] \\
\bottomrule
\end{tabular}
}
\end{table}

